\section{Introducción}
\label{ch:introduccion}

\subsection{Contexto y planteamiento del problema}

La distribución territorial de la contratación pública en Colombia ha sido objeto de creciente atención académica y de política pública. La pregunta central de este proyecto es:

\begin{quote}
\textbf{¿Cómo se distribuye territorialmente la contratación pública en relación con las condiciones socioeconómicas de las unidades territoriales colombianas?}
\end{quote}

Esta interrogante conecta dos dimensiones fundamentales del desarrollo: la asignación de recursos públicos a través de la contratación estatal y las condiciones estructurales de los territorios donde esos recursos se ejecutan. Una correspondencia sistemática (o su ausencia) tendría implicaciones directas para el diseño de políticas de equidad territorial, la focalización de la inversión pública y los mecanismos de control y transparencia.

\subsection{Justificación}

Los municipios colombianos presentan heterogeneidad en sus condiciones socioeconómicas, infraestructura, capacidad institucional y acceso a mercados. La contratación pública también presenta patrones territoriales que pueden responder a necesidades, capacidades administrativas, localización de proveedores y reglas de asignación. Estudiar su asociación con las condiciones socioeconómicas puede aportar evidencia para:

\begin{itemize}
    \item Evaluar si la contratación pública se concentra en municipios con mayores condiciones de vulnerabilidad (enfoque de equidad) o si reproduce desigualdades preexistentes.
    \item Proporcionar evidencia empírica para el diseño de mecanismos de focalización territorial.
    \item Contribuir a la literatura sobre economía política de la contratación pública y desarrollo territorial en contextos de descentralización.
\end{itemize}

\subsection{Antecedentes del proyecto}

El proyecto se origina en la necesidad de construir una base de evidencia rigurosa, reproducible y metodológicamente transparente. Las etapas precedentes han incluido:

\begin{enumerate}
    \item Revisión del estado del arte sobre vulnerabilidad socioeconómica, medición multidimensional y contratación pública en Colombia.
    \item Selección y auditoría de fuentes primarias: Censo Nacional de Población y Vivienda 2018 (CNPV 2018), Índice de Pobreza Multidimensional 2018 (IPM 2018) y DIVIPOLA.
    \item Definición de la unidad de análisis (municipio) y homologación territorial entre fuentes.
    \item Construcción de indicadores socioeconómicos municipales comparables.
    \item Diagnósticos estadísticos preliminares de aptitud para análisis multivariado.
\end{enumerate}

El presente informe consolida metodológicamente estas etapas. Su alcance empírico termina en la preparación y diagnóstico preliminar de indicadores censales de 2018; todavía no contiene resultados sobre contratación pública ni permite responder la pregunta sustantiva.

\subsection{Objetivo del informe}

Documentar de manera rigurosa, reproducible y trazable \textbf{todo el proceso metodológico realizado hasta el estado actual del proyecto}, desde la selección y auditoría inicial de las fuentes de información hasta la construcción de los indicadores socioeconómicos municipales y los diagnósticos preliminares. Este documento es el \textbf{primer informe metodológico formal del proyecto} y está diseñado para crecer modularmente con las etapas subsiguientes.

\subsection{Alcance y limitaciones declaradas}

\begin{itemize}
    \item \textbf{No se afirma} que los indicadores hayan sido aprobados definitivamente para PCA.
    \item \textbf{No se afirma} que se haya ejecutado un Análisis de Componentes Principales (PCA).
    \item \textbf{No se afirma} que exista un Índice de Vulnerabilidad Socioeconómica (IVS) definitivo.
    \item \textbf{No se afirma} que se haya seleccionado definitivamente la estructura factorial.
    \item \textbf{No se afirma} que se haya realizado validación estadística definitiva del índice.
    \item \textbf{No se afirma} que SECOP II haya sido integrado.
    \item \textbf{No se afirma} que se haya realizado análisis espacial.
\end{itemize}

Cuando una etapa esté pendiente, se identifica explícitamente en la sección \ref{sec:etapas_pendientes} y en la tabla de estado del proyecto (Sección \ref{sec:estado_proyecto}).

\subsection{Estructura del documento}

El informe presenta, en este orden, el marco conceptual, el diseño y alcance, las fuentes y unidades territoriales, la selección y tratamiento de datos, la construcción y control de indicadores, los diagnósticos preliminares, el plan de PCA y la reproducibilidad. Cierra con las limitaciones, el estado de avance y las decisiones pendientes.

